\documentclass[a4paper]{article}

% Basic packages
\usepackage[fontset=mac]{ctex}
\usepackage{amsmath, amssymb}
\usepackage{graphicx}
\usepackage[margin=2.5cm]{geometry}
\usepackage[most]{tcolorbox}
\usepackage{etoolbox}
\usepackage{listings}
\usepackage{hyperref}
\usepackage{booktabs}
\usepackage{subcaption}
\usepackage{float}
\usepackage{tikz}
\usetikzlibrary{positioning}
\IfFileExists{pgfplots.sty}{
    \usepackage{pgfplots}
    \pgfplotsset{compat=1.18}
}{}

% Highlight boxes
\newtcolorbox{knowledgebox}[1]{
    enhanced, colback=blue!5!white, colframe=blue!75!black, colbacktitle=blue!75!black,
    coltitle=white, fonttitle=\bfseries, title=#1, attach boxed title to top left={yshift=-2mm, xshift=2mm},
    boxrule=1pt, sharp corners
}

\newtcolorbox{importantbox}[1]{
    enhanced, colback=yellow!10!white, colframe=yellow!80!black, colbacktitle=yellow!80!black,
    coltitle=black, fonttitle=\bfseries, title=#1, sharp corners
}

\newtcolorbox{warningbox}[1]{
    enhanced, colback=red!5!white, colframe=red!75!black, colbacktitle=red!75!black,
    coltitle=white, fonttitle=\bfseries, title=#1, sharp corners
}

% Listings
\lstset{
    language=Python,
    basicstyle=\ttfamily\small,
    keywordstyle=\color{blue},
    stringstyle=\color{red!60!black},
    commentstyle=\color{green!60!black},
    breaklines=true,
    frame=single,
    numbers=left,
    numberstyle=\tiny\color{gray},
    captionpos=b,
    extendedchars=true,
    inputencoding=utf8
}

% Front-page metadata
\newcommand{\notetitle}{翁家翌：OpenAI，GPT，强化学习，Infra，后训练，天授，tuixue，开源，CMU，清华｜WhynotTV Podcast \#4}
\newcommand{\noteauthors}{WhynotTV \& 笔记整理}
\newcommand{\notedate}{\today}
\newcommand{\videochannel}{WhynotTV Podcast}
\newcommand{\videopublishdate}{2025年}
\newcommand{\videoduration}{2小时2分45秒}
\newcommand{\videourl}{https://www.bilibili.com/video/BV1darmBcE4A/}
\newcommand{\videocoverpath}{cover.jpg}

\begin{document}

\begin{titlepage}
\centering
{\Large 播客笔记\par}
\vspace{1.2cm}
{\huge\bfseries \notetitle\par}
\vspace{0.8cm}
{\large \noteauthors\par}
\vspace{0.3cm}
{\large \notedate\par}
\vspace{1.2cm}

\ifx\videocoverpath\empty
{\small 请在生成笔记时填入视频封面图的本地路径。\par}
\else
\includegraphics[width=0.82\textwidth,height=0.45\textheight,keepaspectratio]{\videocoverpath}\par
\fi

\vfill
\begin{tcolorbox}[width=0.9\textwidth, colback=black!2!white, colframe=black!60, sharp corners]
\textbf{视频作者/频道}：\videochannel\par
\textbf{发布日期}：\videopublishdate\par
\textbf{视频时长}：\videoduration\par
\textbf{视频链接}：
\ifx\videourl\empty
未填写
\else
\href{\videourl}{\nolinkurl{\videourl}}
\fi
\end{tcolorbox}
\end{titlepage}

\tableofcontents
\newpage

%% --- 正文内容开始 --- %%

\section{引言：走进翁家翌的世界}

本期播客嘉宾翁家翌，2022年加入 OpenAI，是 GPT-3.5、GPT-4、GPT-4o 到 GPT-5 一系列核心模型背后的核心贡献者之一。他的核心贡献可以总结为三个关键词：\textbf{强化学习（RL）、Post-Training、Infra}。他搭了整个 OpenAI 的 Post-Training RL Infra，因此几乎所有 model release 上都有他的名字。

\begin{figure}[H]
\centering
\includegraphics[width=0.7\textwidth]{figures/talk_opening.jpg}
\caption{播客开场画面 \protect\footnotemark}
\end{figure}
\footnotetext{视频画面时间区间：00:01:01--00:01:12。}

在加入 OpenAI 之前，翁家翌就已经用开源和产品影响过无数人：从清华开源作业与资料试图打破信息差，到开源强化学习框架\textbf{天授（Tianshou）}，再到免费签证查询系统\textbf{tuixue.online}。他把代码和工具视作一种\textbf{慈善}——追求的不是掌声，而是 impact。

\begin{importantbox}{翁家翌的核心画像}

\begin{itemize}
\item \textbf{角色}：OpenAI 核心研究员，Post-Training RL Infra 负责人
\item \textbf{关键词}：强化学习、Post-Training、Infra、开源、Impact
\item \textbf{代表性项目}：天授（Tianshou）RL 框架、tuixue.online 签证查询系统
\item \textbf{教育背景}：清华本科 → CMU 硕士
\item \textbf{核心哲学}：打破信息差、追求 Impact 而非声望、工程能力 > 研究能力
\end{itemize}
\end{importantbox}


\subsection{本章小结}
翁家翌是一个从清华起步、经 CMU 再到 OpenAI 的典型技术人才，但他对"开源"、"打破信息差"、"追求Impact"有着超乎常人的执着。他的故事横跨学术、工业、开源社区，为我们提供了一个观察顶尖 AI 公司和 AI 人才成长的独特窗口。

\section{成长与求学之路}

\subsection{福州的天才少年}

翁家翌的成长故事从福州开始。他从小学一年级开始学奥数，一直学到六年级、初中，发现自己做数学题比谁都快——用现在时髦的话说就是"System 1"，表层意识直接过，看一眼出答案。

但他坦言自己学新东西"比别人慢两到三倍"：

\begin{quote}
读一段代码，我要比别人花很多时间去理解整个 context。但我一旦理解之后，我用得就非常快。就像小时候背书——睡前磕磕巴巴背出来，睡一觉第二天倒背如流。
\end{quote}

这种"慢学快用"的模式贯穿他整个职业生涯。

\begin{warningbox}{学习速度 vs 理解深度}

翁家翌的学习方式是一个重要提醒：理解慢不一定是劣势。深度理解后的"快"往往远超表面的"快"。在 AI 领域尤其是 Infra 工作中，对 context 的完整掌握比快速上手重要得多。
\end{warningbox}


\subsection{高中竞赛与升学的岔路口}

高中时期，翁家翌参加了信息学竞赛。竞赛训练了他对算法的直觉，但他在高三的某一天，突然产生了一个影响他一生的"顿悟"：

\begin{quote}
如果人生是一场游戏的话，你的游戏结算分数是记得你名字的人的数量。
\end{quote}

这个想法被他称为"未来向过去发送信息"——从那一刻起，"被更多人记住"成为他衡量人生价值的核心标准。

\subsection{清华岁月：在评价体系之内与之外}

进入清华后，翁家翌一方面按规则在体系内"做事"——拿 GPA、发 paper、打比赛；另一方面却始终对这个体系保持着清醒的疏离。他的策略是：

\begin{quote}
够用就行。算一下你现在已经多少分了，然后安排考试到底要不要复习。87分（B+）就很满足了，多一分都不想花时间。
\end{quote}

这种"精准投入"的态度，使他把大量精力投入到自己真正在意的事情上——开源、做工具、打破信息差。

\begin{knowledgebox}{评价体系的二重性}

翁家翌区分了两种"外部评价"：
\begin{itemize}
\item \textbf{既有评价体系}（GPA、高考、PhD > Master）：他不喜欢，因为它是一套僵化的标准，推着人走。
\item \textbf{共识评价}（GitHub star、点击量、citation、inference 次数）：他认可，因为这是"每个人发自内心的点赞"，代表真正的 impact。
\end{itemize}
他既不盲从体系，也不完全否定外部评价——他要的是"发自内心的认可"而非"体制的盖章"。
\end{knowledgebox}


\subsection{留学申请受挫与挣脱}

2020年疫情期间，翁家翌申请了18所美国学校，只收到 Google 和 AutoML 两个 offer。这对于一个清华绩点不低、有 paper 有开源项目的人来说是巨大的打击。但他选择把精力投入到更有意义的事情上——在家写天授、写 tuixue.online。

\begin{quote}
转移注意力，不要天天去关注宏大的国际叙事，而是专注于手头上的事情——这样可能让自己内心更平静一些。
\end{quote}

最终他去了 CMU 读硕士，这也成为他进入 OpenAI 的关键跳板。

\subsection{本章小结}
翁家翌的成长轨迹揭示了一个核心张力：如何在既有评价体系中保持自己的内在标准。他的策略是"够用就行"——在体系内满足最低门槛，然后把全部精力投入到自己真正认为有价值的事情上。这种态度的底层是对"impact"的执着追求，而非对"声望"的盲目追逐。

\section{开源精神与 Impact 哲学}

\subsection{清华开源作业：打破信息差的第一枪}

在清华期间，翁家翌做了一件在当时校园里被视为"破坏规则"的事——把课程作业和资料全部开源。

\begin{quote}
信息差是一个如果你在清华生存的话是很有用的东西。但我觉得每个人都应该平等地拥有这个信息。哪怕别人有了同样的认知也做不出来，我还是觉得应该打破它。
\end{quote}

这种想法在当时看来有些天真，但十几年后回头看，它其实是对"知识应该流动"这个价值观的朴素坚持。

\begin{knowledgebox}{开源的意义}

翁家翌对开源的看法不是"技术炫耀"或"社区贡献"，而是一个更根本的信念：\textbf{知识不应该成为特权的护城河}。他的开源行为本质上是对"信息差"这种社会不平等的对抗——哪怕这种对抗在当时来看是"自毁长城"。
\end{knowledgebox}


\subsection{天授（Tianshou）：两周写完的 RL 框架}

2020年初，翁家翌想把之前写的 RL 实验代码整合成一个框架。他先看了当时最流行的 RLlib，花了一个月——

\begin{quote}
太复杂了，抽象太多了。一个 RLlib 快几十万行代码，完全不被人接受。我根本不知道我要改的东西该怎么实现。
\end{quote}

于是他决定推倒重来。\textbf{两周}，一个人，写出了天授的第一版。

\begin{importantbox}{天授成功的关键}

\begin{enumerate}
\item \textbf{抓住用户需求}：Researcher 需要"好用、好改"的 RL 框架，天授做到了。
\item \textbf{一致性优先}：一个人写，顶层设计一致，代码短，抽象清晰。
\item \textbf{修改点唯一}：如果要在天授上加 feature，只有一个地方能改——设计本身就规定了修改路径。
\item \textbf{应用于心}：每个算法实现不超过20行，对着 paper 实现就好。
\end{enumerate}
\end{importantbox}


他对天授的定义是"做慈善"——完全 nonprofit，不期望任何回报。这种态度和他的 impact 哲学一脉相承：

\begin{quote}
做慈善项目让我感觉非常满足。
\end{quote}

\begin{warningbox}{项目腐化的根源：不一致性}

翁家翌指出了大多数开源项目腐化的核心原因：\textbf{多人在不了解彼此假设的情况下各自写代码}，导致复制粘贴、代码膨胀、抽象泄漏。一致性比代码行数、测试覆盖率等任何工程指标都更重要。好项目从头到尾都是一致的。
\end{warningbox}


\subsection{tuixue.online：签证查询系统}

在疫情最严重的时候，他做了另一个"慈善"项目——\texttt{tuixue.online}，一个免费签证查询系统，帮助无数留学生追踪签证状态。同样是 non-profit，同样是出于"做一些对大家有意义的事"的冲动。

\subsection{Impact 的底层逻辑}

当被问到"为什么不自己过得开心就好"时，翁家翌给出了一个存在主义式的回答：

\begin{quote}
人生是一种体验。既然你已经来到这个世界上了，那就不要浪费这段旅程。
\end{quote}

他区分了两种"被看见"：
\begin{itemize}
\item \textbf{名望（fame）}：坐上了某个位置，有头衔有身份，但对别人可能毫无意义甚至负面。
\item \textbf{Impact}：做了对别人有用的事，被真正受益的人记住。
\end{itemize}

\begin{quote}
我想做的是力所能及地对身边的人好，做一些对大家有意义的事。这个度量标准是：让更多的人记得你。
\end{quote}

但他也清醒地认识到：\textbf{这个标准只对自己适用}，不会拿它去要求别人。"如果你发现自己的标准成了新的枷锁——你可以改。你不是这个标准的奴隶。"

\subsection{本章小结}
翁家翌的开源行为不是社区贡献策略，而是价值观的表达：打破信息差、降低门槛、做对别人有用的事。他的 impact 哲学既反对既有评价体系的僵化，也警惕自己的标准变成新的"成规"。天授和 tuixue.online 是他这一哲学最完美的注脚。

\section{加入 OpenAI}

\subsection{CMU 硕士：用工程能力差异化竞争}

翁家翌选择读硕士而非 PhD，理由非常务实：

\begin{quote}
如果你想进工业界，读 PhD 就是浪费生命。你完全可以用 master 为跳板，在本科或硕士期间攒够 citation，做出与众不同的项目，让自己可以和 PhD candidate 同台竞技。
\end{quote}

他投了18家公司，最初只有 Google 和 AutoML 两个 offer。意识到需要换个策略后，他重新面试，先后拿到了幻方（后来的 DeepSeek）、英伟达和 TikTok 的 offer——以及 OpenAI。

\begin{knowledgebox}{为什么放弃 PhD？}

翁家翌的核心判断：
\begin{itemize}
\item \textbf{PhD 培养的是学术能力}：写 paper、讲故事、画图、宣发。
\item \textbf{工业界需要的是工程能力}：建 infra、快速验证 idea、保证正确性。
\item \textbf{差异化竞争}：让公司愿意招"master 的你"而不是"另一个 PhD"——你必须有 PhD 候选人所不具备的东西。
\end{itemize}
\end{knowledgebox}


\subsection{与 John Schulman 的面试}

翁家翌加入 OpenAI 的面试由 John Schulman 亲自进行。Schulman 出了一道非常 open-ended 的题，给三小时。翁家翌两小时做完，虽然展示时出了个 bug，但当场修好。

Schulman 为什么招他？

\begin{quote}
他觉得我 GitHub 非常漂亮。招一个有良好工程能力的人进来，对任何项目都是有益的。
\end{quote}

翁家翌在 Schulman 离职那天难过了一个下午，"把电脑关了，什么都不干"。

\subsection{为什么选择 OpenAI（在 ChatGPT 之前）}

2022年加入 OpenAI 时，ChatGPT 还没有发布。翁家翌的选择基于两点：

\begin{enumerate}
\item OpenAI 和 DeepMind 是当时 RL 领域做得最好的两个 lab，能进去就是荣耀。
\item 他想了解"世界上最前沿的 research 到底是怎么做的"——不是学校里的"小作坊"模式，而是有方法论的工业级研究。
\end{enumerate}

\subsection{本章小结}
从 CMU 到 OpenAI，翁家翌走的是一条"工程驱动"的路径。他不是传统的 research PhD，而是一个"用工程能力说话"的人。John Schulman 看中的正是这一点——漂亮的 GitHub 页面比 paper 数量更有说服力。

\section{研究 vs 工程：Infra 是核心竞争力}

\subsection{工程能力为什么比研究能力重要}

翁家翌引用了一位同事的话：

\begin{importantbox}{核心命题}{
\textbf{教一个 researcher 如何做好 engineering，要远比教一个 engineer 如何做好 research 来得难。}
}
\end{importantbox}

他的论证链条非常清晰：

\begin{enumerate}
\item \textbf{Idea 非常便宜}：找人讨论一下就有了。真正有研究直觉的人（如 Alec Radford）动脑子比普通 PhD 动脑子有用得多。
\item \textbf{验证 Idea 才是瓶颈}：你单位时间内能验证多少有效的 idea，并且保证验证结果是正确的。
\item \textbf{验证依赖 Infra}：Infra 有 bug → 验证结果不可靠 → 研究方向被误导。
\item \textbf{谁修 bug 谁赢}：\textbf{每家的 Infra 都有不同程度的 bug，谁修 bug 越多，谁的模型训得越好。}
\end{enumerate}

\begin{quote}
如果你把 bug 全修了，那有可能算法连改都不用改——就很好。
\end{quote}

\subsection{Infra 的"迭代速度"是生死线}

翁家翌对 AI 公司核心竞争力的定义极其简洁：

\begin{importantbox}{Infra 生死线}

$$ \text{竞争力} \approx \frac{\text{单位时间内正确迭代的次数}}{\text{bug 数量}} $$
\begin{itemize}
\item \textbf{分子}：迭代速度（iteration speed）
\item \textbf{分母}：Infra 正确性（有多少 bug 没修）
\end{itemize}
\end{importantbox}


他进一步指出，目前所有大模型公司的 Infra"都有 bug"，差距只是多少的问题。而"修 bug"这件事不能外包给纯 researcher——他们"不愿意修也不想修，因为这不在他们的评价体系里"。

\begin{knowledgebox}{Infra 工程师 vs Researcher 的分工}

\begin{itemize}
\item \textbf{Researcher}：找人讨论 → 产生 idea → 改 flag 跑实验 → 看结果。不需要了解 Infra 底层。
\item \textbf{Infra Engineer}：设计 distributed training 架构、修 bug、保持一致性、优化吞吐量。
\item \textbf{正确的关系}：Researcher 提需求，Infra Engineer 建底座。底座越稳，researcher 越能专注于 idea 验证。
\end{itemize}
\end{knowledgebox}


\subsection{本章小结}
翁家翌的"Infra 至上论"是对当前 AI 研究范式的一剂清醒剂。在 scaling law 趋于平缓的今天，谁能更快、更正确地验证 idea，谁就能在竞争中领先。而这背后，Infra 的质量和迭代速度是决定性的——"算法改不改都行，先把 bug 修完"。

\section{RLHF 与 Post-Training 技术深度}

\subsection{什么是 Post-Training}

当被问到"什么是强化学习和 post-training"时，翁家翌给出了一个简洁的定义：

\begin{quote}
Post-training 就是模型在预训练（pre-training）完成之后的进一步训练阶段。RLHF（Reinforcement Learning from Human Feedback）是 post-training 的核心方法之一——用人类反馈作为奖励信号，通过强化学习来对齐模型的行为。
\end{quote}

\subsection{2022年 RLHF 的关键挑战}

回到2022年（ChatGPT 发布前夕），RLHF 面临的核心挑战可以归纳为：

\begin{enumerate}
\item \textbf{Infra 稳定性}：工业级 RL 训练涉及多个模型（policy model、reward model、reference model），分布式训练的一致性和容错是巨大挑战。
\item \textbf{规模扩展}：从实验室的小规模实验扩展到 GPT 级别的模型，每一层抽象都可能断裂。
\item \textbf{反馈质量}：人类标注的质量和一致性直接影响 RLHF 效果——"垃圾 reward 只会训练出垃圾 policy"。
\end{enumerate}

翁家翌的贡献正是搭建了这一整套 Post-Training RL Infra，使得从 GPT-3.5 开始的每一次模型发布都有了可靠的 RLHF pipeline。

\subsection{工业级 RL Infra 的挑战}

他将工业级 RL Infra 的挑战归纳为几个核心维度：

\begin{itemize}
\item \textbf{吞吐量}：单位时间能处理多少训练样本和奖励信号。
\item \textbf{正确性验证}：如何确保分布式训练中每一步的结果都是可靠的。
\item \textbf{Context 一致性}："人多了，context inconsistent，每个人都会写一些奇怪的东西"——这是大型 Infra 项目的自然熵增。
\item \textbf{迭代周期}：从改代码到看到实验结果的时间间隔越短越好。
\end{itemize}

\begin{figure}[H]
\centering
\begin{tikzpicture}[node distance=0.8cm, auto]
  \tikzstyle{layer}=[rectangle, draw, minimum width=12cm, minimum height=0.8cm,
                     align=center, font=\small, rounded corners=1pt]
  \node [layer, fill=blue!10] (exp)    {实验层：Researcher 提 idea、设 reward 函数、改 flag 跑实验};
  \node [layer, fill=green!10, below of=exp] (sched) {调度层：分布式训练任务编排、资源分配、容错恢复};
  \node [layer, fill=yellow!10, below of=sched] (rl) {RL 引擎：Policy Model + Reward Model + Reference Model 协同训练};
  \node [layer, fill=orange!10, below of=rl] (data) {数据层：人类反馈收集、标注质量控制、Reward Model 训练};
  \node [layer, fill=red!10, below of=data] (infra) {基础设施层：GPU 集群、网络、存储、分布式通信};
\end{tikzpicture}
\caption{Post-Training RL Infra 分层架构}
\end{figure}

\subsection{Agent RL 与标准 RL 的关系}

关于当前热门的 Agent RL：

\begin{quote}
没有本质差别。本身就是同一个东西，中间多加了几个 token。环境方面的改变而已，谈不上新的挑战。
\end{quote}

这个判断来自一线实践者，值得深思——很多看起来"新"的东西，本质上还是老问题的变形。

\subsection{本章小结}
RLHF/Post-Training 的核心挑战不在算法创新，而在 Infra 工程。"搭 Infra"这项工作决定了整个 RLHF pipeline 能否正确、高效地运转。翁家翌的观点是：在大模型时代，Infra 工程能力就是最硬核的研究能力。

\section{OpenAI 内部视角}

\subsection{2022年初印象：人才密度与组织架构}

2022年加入 OpenAI 时，翁家翌的第一印象是——

\begin{quote}
人才密度极高。每个人都在做自己最擅长的事情，各司其职。组织架构非常扁平，信息流动快。
\end{quote}

但随着公司规模扩大，组织问题也开始浮现。"人多了，context inconsistent，每个人都会写一些奇怪的东西"——这不仅是代码层面的问题，更是组织协作的根本挑战。

\subsection{Sam Altman 被开除事件：内部视角}

2023年11月，OpenAI 董事会突然宣布开除 CEO Sam Altman。外部传言纷飞——"Ilya 到底看到了什么？"

\begin{importantbox}{翁家翌的内部视角}

\begin{itemize}
\item \textbf{原因}：Ilya 和其他董事会成员不信任 Sam，投票把他投出去了。
\item \textbf{底层员工的反应}：非常震惊（surprise），完全不知道发生了什么。
\item \textbf{董事会缺乏透明度}：董事会之前对底下的人是缺乏透明度的，决策过程不透明。
\item \textbf{Sam 为什么能回来}：因为很多员工觉得"如果纯技术出身的人领导，可能没有那么多资源"——AGI 的实现不只是技术问题，还需要融资、商业、资源整合。
\end{itemize}
\end{importantbox}

\begin{quote}
你可以把 Sam 抽象成一个 Persona，这个 Identity 乘以一个系数。别人对这个 Identity 的认同感如果缺失了，组织就运转不起来。这个短时间内 AI 替代不了。
\end{quote}

\subsection{人才流失：一个健康的组织没有人是不可替代的}

面对 OpenAI 持续的人才流失（包括 John Schulman、Ilya Sutskever 等人的离开），翁家翌的态度非常冷静：

\begin{quote}
一个健康的组织是所有人都是可以替代的。你只要能持续培养新人、有造血能力、让组织正常运转就行。虽然走了很多人，但可以花时间精力再培养一波新人——像干细胞一样。
\end{quote}

这意味着：\textbf{OpenAI 今天做的东西，外面的人也不是说难如登天、不能复刻。} 因为没有黑魔法——"把最简单的东西做好就好了"。

\subsection{AI 竞赛的生死线：迭代速度}

当被问到 OpenAI 面对 DeepSeek 等竞争对手的压力时，翁家翌指出：

\begin{importantbox}{真正让 OpenAI 警觉的事}{
不是 DeepSeek 的模型在哪个榜单上比 GPT 分高——OpenAI "很长时间以来都没有特地为了刷 benchmark 做什么事情"。\textbf{真正让他们警觉的是 DeepSeek 的迭代速度}——也就是 Infra 的 cycle time。
}
\end{importantbox}

但他也坦言：OpenAI 目前在这个指标上\textbf{不是世界第一}。

\begin{quote}
你抽出一个团队去搞 startup，他们的迭代速度斜率肯定远比 OpenAI 高。代码库小、沟通成本少、只集中在一个 use case。但 OpenAI 要同时考虑很多 use case——这是组织大了之后的必然问题。每个公司做大了都会变慢。最后比的不是谁更好，而是谁"不是那么差"。
\end{quote}

\begin{knowledgebox}{大公司 vs 小团队的速度悖论}

\begin{itemize}
\item \textbf{小团队优势}：代码库小、context 一致、沟通成本低 → 迭代快
\item \textbf{大公司优势}：用户反馈多、资源充裕、多 use case 测试 → 鲁棒性强
\item \textbf{翁家翌的判断}：这个 trade-off 是"人类组织发展到一定规模后必然会面对的问题，无法避免"
\item \textbf{一个可能的解法}：拥有无限 context 的 AI Agent 来负责组织的信息共享和决策——替代人类 CEO
\end{itemize}
\end{knowledgebox}


\subsection{OpenAI 还 Open 吗？}

对于"OpenAI 已经和 Open 没什么关系了"的批评，翁家翌承认这里存在一个 trade-off：

\begin{quote}
你没有办法直接把最好的模型开源——因为公司要生存。公司如果不能生存，就没有办法继续融资、做实验、有突破性进展。这都是很现实的问题。
\end{quote}

他将 OpenAI 的使命拆成两步：\textbf{先实现 AGI，再造福全人类}。而"造福全人类"目前的理解是做产品、以尽可能便宜的价格让普通人接触到技术（如免费的 ChatGPT），而不是直接丢一个裸权重出去。

他还透露了一个有趣的细节：John Schulman 曾问过他"要不要把 RLHF 的内容开源"，他当时的回答是"不太好吧，为了公司的考量"。

\begin{warningbox}{OpenAI 的闭源是博弈论结果}

翁家翌的视角：即使 OpenAI 想开源，如果"有人就是想挣钱"，开源就等于把技术送给竞争对手闭源商用。为了防止这种情况，OpenAI "不得不闭源"。本质上是一个囚徒困境——因为不是所有人都同一条心。
\end{warningbox}


\subsection{本章小结}
从内部看 OpenAI，最核心的洞察是：AGI 的实现不只是技术问题，更是组织问题、博弈论问题、资源整合问题。Sam Altman 的价值在于他能在商业和资源层面让组织运转起来，而 Infra 的迭代速度是技术层面的生死线。人才流失不可怕——可怕的是失去造血能力。

\section{未来展望：AGI、宿命论与十年愿景}

\subsection{AGI 的定义与距离}

OpenAI 内部有一个笑话："你抓15个人，可能有20种定义 AGI 的方法。"

翁家翌的定义是：

\begin{quote}
如果这个东西能完成我80-90\%认为有意义的 task，那可能就是 AGI 了。
\end{quote}

按照这个标准，他认为 AGI 还没有实现。核心障碍是：AGI 需要在极度 out-of-distribution 的场景下工作——比如修改 Infra 代码。"AI Infra 的代码在训练数据里的占比几乎为零"，而验证反馈的时间又太长，成本太高。

\begin{quote}
Strawberry 出来之前，内部已经用了一段时间。大家都觉得"我的工作要被取代了"，甚至有人开玩笑说"写一堆屎山吧，后面 Strawberry 会帮我们清理的"。一两年过去了，屎山还在那儿——并不会真正改变什么东西。每个人都会 over-react，觉得技术来了就天翻地覆，但实际上它是一个很慢很慢的循序渐进的过程。
\end{quote}

\subsection{宿命论：世界是确定性的吗？}

在播客的最后阶段，翁家翌抛出了一个极其"哲学"的观点：

\begin{importantbox}{翁家翌的宿命论}

两个核心主张：
\begin{enumerate}
\item \textbf{世界是确定性的}：所有东西都可以被预测。你现在脑子里在想什么，下一个单词说什么，全都是宇宙大爆炸那一刻就定好了——这是一个确定性的马尔可夫过程。
\item \textbf{人是没有自由意志的}：他声称这是"验证过无数遍"的结论。
\end{enumerate}
\end{importantbox}


当被问到如果有能预测未来的 AI model，最好的选择是什么时——

\begin{quote}
毁掉它。让这样的 AI model 永远不要出来。因为它会导致所有价值体系的崩塌。
\end{quote}

但他又补充："有些人就非常愿意去开发这种 model——不是为了逃脱宿命，而是为了搞清楚世界背后运行的规律，搞清楚为什么世界是宿命论的。"

他的应对策略：

\begin{quote}
最好的方式就是忘掉这一切，假装你不知道这个事，然后去体验当前的经历。我一直都是这么做的。
\end{quote}

\subsection{十年愿景}

被问到希望十年后的自己是什么样时，翁家翌的回答回归到了他高中时的那个顿悟：

\begin{quote}
我希望十年后的我，还是能做一些对大家有意义的事。让更多的人记得我。
\end{quote}

他提到最想用 AI 解决的难题是\textbf{预测未来}——"我个人对于造一个世界还是很有吸引力的。你需要提前生成剧本。"

\subsection{本章小结}
从技术维度到哲学维度，翁家翌的思考呈现出一种罕见的深度。他既是"Infra 至上"的极致务实派（修 bug、提吞吐），又是一个宿命论者（世界是确定性的马尔可夫过程）。这种张力恰恰反映了他作为 AI 研究者的复杂面向：在最底层的工程细节和最顶层的人生意义之间自由切换。

\section{总结与延伸}

\subsection{讲者的核心观点回顾}

翁家翌在播客结尾没有给出传统的"总结"，但他的整场谈话中反复出现的主题可以归纳为：

\begin{enumerate}
\item \textbf{Impact > Fame}：追求实质影响而非空名。
\item \textbf{工程 > 研究}：在 AI 当前阶段，工程能力（建 Infra、修 bug、提吞吐）比研究能力（发 paper、写故事）更重要。
\item \textbf{迭代速度 = 核心竞争力}：谁的 Infra 迭代快、bug 少，谁的模型就训得好。
\item \textbf{信息差应该被打破}：知识不应成为特权护城河。
\item \textbf{世界可能是确定性的}：但最佳策略是"忘掉它，去体验"。
\end{enumerate}

\subsection{结构化蒸馏}

通过这场播客，我们可以提取出几个核心洞察：

\begin{importantbox}{翁家翌式 AI 人才观}

\begin{itemize}
\item \textbf{差异化定位}：Master 要和 PhD 竞争，靠的不是"更像 PhD"，而是"PhD 做不到的事"——工程能力。
\item \textbf{够用就行}：在评价体系中满足最低门槛，把剩余精力投入自己真正在意的方向。
\item \textbf{开源是价值观}：不是社区运营策略，而是对信息不对称的对抗。
\item \textbf{Infra 是硬通货}：在这个阶段，能修 bug 的人比能提 idea 的人更稀缺。
\end{itemize}
\end{importantbox}


\begin{importantbox}{翁家翌式 AI 竞赛观}

\begin{itemize}
\item \textbf{不追 benchmark}：OpenAI 不关心谁在某榜上比 GPT 高多少分。
\item \textbf{追迭代速度}：真正让 OpenAI 警觉的是 DeepSeek 的 Infra 迭代效率。
\item \textbf{大公司的宿命}：规模带来 context 不一致和臃肿——这是所有组织发展到一定阶段的必然结果。
\item \textbf{出路可能是 AI Agent}：一个拥有无限 context 的 agent 来解决人类组织的协调问题。
\end{itemize}
\end{importantbox}


\subsection{扩展思考}

翁家翌的对话提出了几个值得深思的开放问题：

\begin{enumerate}
\item \textbf{Infra 工程师 vs Researcher 的分工是否可持续？} 当 Infra 越来越复杂，Researcher 对底层的了解越来越浅，两者之间的 gap 会不会导致研究方向被 Infra 的局限性所绑架？
\item \textbf{"评价体系"的重构}：翁家翌反对 GPA/paper 这类既有评价体系，但提倡 GitHub star/用户量这类"共识评价"。问题是：当共识评价也成为主流后，会不会同样腐化？他自己的回答是"你可以改自己的标准"——但这需要持续的自我觉察。
\item \textbf{开源 vs 闭源的博弈论困境}：OpenAI 的闭源策略在当前市场环境下是理性的，但长期来看会不会削弱整个 AI 生态的进步速度？DeepSeek 的 open weights 是否会改变这个博弈格局？
\item \textbf{AGI 的"out-of-distribution"难题}：翁家翌认为 AGI 的瓶颈不在于已有能力的泛化，而在于处理训练数据中占比极低的领域（如 Infra 代码）。这暗示了一个 deeper problem：AGI 需要的不是更好的 scaling，而是某种根本性的架构创新。
\end{enumerate}

\subsection{拓展阅读}

\begin{itemize}
\item 天授（Tianshou）RL 框架：\href{https://github.com/thu-ml/tianshou}{github.com/thu-ml/tianshou}
\item tuixue.online 签证查询系统：\href{https://tuixue.online}{tuixue.online}
\item OpenAI Post-Training 相关技术博客：\href{https://openai.com/research}{openai.com/research}
\item John Schulman 关于 RLHF 的演讲与论文
\item DeepSeek 技术报告：了解竞争对手的 Infra 设计与迭代效率
\end{itemize}

\end{document}